\section{Naples, the Silence, and Fossanova}

\subsection{The last teaching year}

At Naples in 1272--73 Thomas organized the province's theology school,
lectured, and pressed on with the Third Part of the \work{Summa}---God
incarnate, the Passion, the sacraments, reaching the Eucharist and
penance. In Lent 1273 he did something the professional theologian was
not required to do: he preached to the people of Naples, in their own
tongue, a course of instructions on the Creed, the Our Father, the Ten
Commandments, and the Hail Mary---catecheses whose plan, as
Benedict~XVI observed in presenting them, is recognizably the plan of
the Church's present Catechism: what Christians believe, pray, and do.
To these Naples months the canonization dossier attaches the most famous
of the visionary stories, the crucifix in the chapel of St.~Nicholas
that spoke---``Thou hast written well of me, Thomas; what reward wilt
thou have?'' ``None other than Thyself, Lord''---examined with the other
traditions in section~\ref{sec:traditions}.

\subsection{6 December 1273}

On the morning of St.~Nicholas's day, 6 December 1273, Thomas celebrated
Mass, and then stopped. He wrote no more, dictated no more; the
\work{Summa} halts in the questions on penance. To his companion
Reginald of Piperno, who pressed him to return to work, he gave the
answer the whole world knows in the form the old Encyclopedia renders:
``I can do no more. Such secrets have been revealed to me that all I
have written now appears to be of little value''---or, in the wording
modern accounts prefer, that all of it seemed to him ``like straw.''

\witnesslimit{The witness chain deserves to be stated exactly, because
the remark has become the most quoted sentence Thomas never wrote down.
The event---the cessation itself---is secure: the unfinished
\work{Summa} is its permanent evidence, and every strand of the
tradition presupposes it. The words rest on one deposition: Bartholomew
of Capua, high official of the kingdom of Naples, testifying at the
canonization inquiry of 1319, forty-six years later---and not as an
eyewitness. He recounted what had come to him through Dominican
intermediaries from Reginald, to whom the words were spoken (the
standard modern reconstructions print the chain; this study verified it
at reported level, the process text itself remaining uninspected). The
English renderings differ---``straw,'' ``of little value''---because
they translate a second-hand Latin memory. What the historian may
assert: Thomas stopped; his intimates connected the stopping with an
overwhelming experience at Mass; the ``straw'' formulation is how the
canonization generation remembered his own explanation. What no one can
verify: the sentence's exact wording, or the inner event it gestures
at. This study certifies neither a mystical grace nor a medical
collapse; the sources permit both vocabularies and settle neither.}

\subsection{The road to Lyon and the death at Fossanova}

Gregory~X had convoked a general council at Lyon for 1 May 1274, to
attempt reunion with the Greek church, and ordered Thomas to attend,
bringing his \work{Contra errores Graecorum} (so the 1912 Encyclopedia).
Obedient, the sick man set out in January 1274 with Reginald. On the
road he failed: the canonization-era memory tells of a blow to the head
against a fallen tree, then collapse near Terracina; he was carried to
the castle of Maenza, home of his niece Francesca, and, sensing the end,
asked to die in a religious house. The nearest was the Cistercian abbey
of Fossanova. There, the tradition continues, he dictated at the monks'
request a brief commentary on the Song of Songs---if so, his last
teaching, and it does not survive; there, before receiving the Viaticum,
he made the profession the dossier preserves, submitting all he had
written ``to the judgment and correction of the Holy Roman Church, in
whose obedience I now pass from this life.'' He died in the early hours
of 7 March 1274, aged about forty-nine. The date is as secure as any in
his life; the deathbed speeches are received tradition, recorded here
as such. The Cistercians buried him before their high altar, and kept
him---a possession that would itself become a story
(section~\ref{sec:traditions}).
